% !TeX program = lualatex
\RequirePackage{iftex}
\ifLuaTeX\else
  \GenericError{}{This document must be compiled with LuaLaTeX.}{}{Run `latexmk -lualatex sample.tex` or set your editor's TeX engine to LuaLaTeX.}
  \expandafter\endinput
\fi
\documentclass[11pt]{beamer}
\usepackage{luatexja}
\usepackage{mathtools}
\usepackage{amssymb}
\usepackage{amsthm}
\usepackage{braket}
\usepackage{bm}
\usepackage{booktabs}
\usepackage{array}
\usepackage{tikz}
\usetikzlibrary{arrows.meta, calc}

\setbeamertemplate{proof begin}{}
\setbeamertemplate{proof end}{}

\usetheme[block=fill]{metropolis}

\setbeamersize{
  text margin left=5mm,
  text margin right=5mm
}
\newcommand{\widespacing}{\linespread{1.5}\selectfont}

%%%%%%%%%%%%%%%%%%%%%%%%%%%%%%%% １ページ目 %%%%%%%%%%%%%%%%%%%%%%%%%%%%%%%%%
\title{教員採用試験対策}
\author{今泉 力}

\begin{document}

\begin{frame}
  \titlepage
\end{frame}

%%%%%%%%%%%%%%%%%%%%%%%%%%%%%%%%% ２ページ目 %%%%%%%%%%%%%%%%%%%%%%%%%%%%%%%%%
\begin{frame}
  \frametitle{なぜ教員を志望しましたか。}
  {\widespacing
  私が教員を志望した理由は、子どもたちに数学の面白さ、数学がどのように社会に関わっているかを伝えたいと思ったからです。\\
  私は大学で数学や物理を学習する中で、大学の数学と高校の数学に大きな差があることを痛感いたしました。\\
  しかし、その差を痛感したからこそ、高校数学の段階で「大学やもっと先の将来にどう使われるのか」を理解することに大きな意味があると思いました。\\
  そのため、数学と社会のつながりを生徒に伝え、生徒のキャリアをサポートできるような教員になりたいと考え、志望させていただきました。
  }
\end{frame}

%%%%%%%%%%%%%%%%%%%%%%%%%%%%%%%%%% ３ページ目 %%%%%%%%%%%%%%%%%%%%%%%%%%%%%%%%%
\begin{frame}
  \frametitle{なぜ、当自治体を志望しましたか。}
  {\widespacing
  一番の理由は、やはり自分が生まれ育った地元である愛知県で働き、恩返しがしたいという思いが強いからです。\\
  愛知県は「モノづくり」に関する風土が根付いており、私が志望する高校数学は、そのようなあらゆるモノづくりや論理的思考の「土台」となる学問であると考えています。\\
  自分が愛着を持っているこの地元で、数学を通して、子どもたちの将来の可能性を広げるお手伝いができればと思い、志望しました。\\
  }
\end{frame}

%%%%%%%%%%%%%%%%%%%%%%%%%%%%%%%%%%% ４ページ目 %%%%%%%%%%%%%%%%%%%%%%%%%%%%%%%%%
\begin{frame}
\frametitle{なぜこの校種を志望しましたか。}
{\widespacing
私が高校を志望した理由は、小学校や中学校と比べて内容が専門的になる分、社会とのつながりをより伝えることができると考えたからです。\\
高校生は、進学や就職など自分の将来について具体的に考え始める時期だと思います。\\
高校数学の専門性を活かすことで、生徒が将来について考えるきっかけを作りたいと思い、この校種を志望しました。\\
}
\end{frame}

%%%%%%%%%%%%%%%%%%%%%%%%%%%%%%%%%%% ５ページ目 %%%%%%%%%%%%%%%%%%%%%%%%%%%%%%%%%
\begin{frame}
\frametitle{当自治体の求める教師像を知っていますか。共感している点はどこですか。}
{\widespacing
はい。『あいちの教育ビジョン2030』に掲げられている『持続可能な社会の創り手を育む』という理念だと認識しております 。\\
その中で私が特に共感したのは、『数学が将来役に立つと思われていない』という課題です 。愛知県は自動車産業などのモノづくりが盛んで、実はすごく数学が活きる地域だと思うのですが、それが生徒たちに伝わりきっていないことに、少しもどかしさを感じていました。\\
社会につなげる『キャリア教育』という視点でも、数学を通して子供たちの成長をサポートしていきたいと考えています。\\
}
\end{frame}

%%%%%%%%%%%%%%%%%%%%%%%%%%%%%%%%%%% ６ページ目 %%%%%%%%%%%%%%%%%%%%%%%%%%%%%%%%%
\begin{frame}
  \frametitle{教員の魅力は何ですか}
  {\widespacing
  教員の魅力は、将来の担い手となる学生たちの成長を近くで見守ることができることだと思います。\\
  社会の中で最も若いともいえる学生の成長に、仕事として最も近くで関わることができるのは教員だと思います。\\
  実際に塾講師として働いているとき、2,3年と長期的に指導してきた学生の成長は、私の想像を超えるものがありました。\\
  はじめは問題に対して手が止まっていた生徒が、少しずつ自分の考えで手を動かすようになったり、「特にやりたいことがない」と言っていた生徒が、「英語を勉強したいから外国語大学に行きたい」と言い出したりと、成長の早い彼らの様子を見守ることは、教員の最大の魅力だと思います。\\
  }
\end{frame}

\end{document}
